\vspace{-0.3cm}\section{Use Cases}
\label{sec:usecase}
We demonstrate the value of LACK with two use cases. 
The first use case identifies entities that may be difficult to trace their connection without integrating the data sources\footnote{We present more of these examples on our website: \url{https://climatesense-project.eu/lack/case-studies.html}.}
The second use case cross matches LACK Collectives with the attendees the attendees
most connected LACK Collectives with the attendance list of COP30 (the United Nations Climate Change Conference), and identifies potential investigations.
\begin{sloppypar}
\paragraph{Case study: LACK enables integration of multiple independent sources.} 
We provide two examples with the same structure: two statements, each drawn from a different source — the DeSmog Climate Disinformation Database and LobbyMap — share a common entity that acts as a bridge. 
Neither source mentions the entities cited in the other — yet their combination, through the shared bridge entity, reveals a relationship that no single source makes visible\footnote{The connections revealed are structural, derived from publicly available data; they do not imply wrongdoing, shared intent, or direct collaboration, and are presented as research leads rather than evidential claims.}.

\textit{Example 1: Darren Woods $\rightarrow$ Exxon $\rightarrow$ Heartland Institute}
\noindent\textbf{Bridge entity:} Exxon. \textbf{(InfluenceMap)} Darren Woods \texttt{lack:leadsAt} Exxon; \textbf{(DeSmog)} Exxon \texttt{lack:sponsored} Heartland Institute.
InfluenceMap documents Woods leading Exxon; DeSmog records Exxon as a sponsor of the Heartland Institute, a think tank prominently associated with climate science denial.
Neither source references the other's entity, yet LACK makes the chain explicit, linking a named corporate officer to a climate change denial organisation through his employer --- a pattern generalised by joining \texttt{lack:leadsAt} with \texttt{lack:sponsored} in a single SPARQL query.

\textit{Example 2: Cambridge Analytica $\rightarrow$ BP $\rightarrow$ Business Council of Australia}
\noindent\textbf{Bridge entity:} BP. \textbf{(DeSmog)} Cambridge Analytica \texttt{lack:hasPartner} BP; \textbf{(InfluenceMap)} BP \texttt{lack:memberOf} Business Council of Australia.
DeSmog records a partnership between the data and political consultancy firm Cambridge Analytica and BP, while InfluenceMap separately documents BP's membership of the Business Council of Australia, an industry body with a record of opposing climate legislation.
The two facts originate from different investigative contexts --- political data operations vs.\ fossil fuel industry lobbying in Australia --- and share only BP, yet LACK surfaces this cross-domain, cross-jurisdictional structural connection that neither source exposes alone.
\end{sloppypar}

Tracing such chains allows researchers to surface accountability relationships that are structurally present in the data but invisible within any single source.

% %-----------------------------------------------------------------------
% \subsection*{Case 3: Nick Stavropoulos $\rightarrow$ PG\&E $\rightarrow$ ZETA}

% \noindent\textbf{Source A (DeSmog):} Nick Stavropoulos \texttt{lack:leadsAt} PG\&E.\\
% \textbf{Source B (LobbyMap):} PG\&E \texttt{lack:memberOf} Zero Emission Transportation Association (ZETA).\\
% \textbf{Bridge entity:} PG\&E.

% \smallskip
% \noindent DeSmog profiles Nick Stavropoulos as an executive at PG\&E, the Californian utility. LobbyMap independently records PG\&E's membership of the Zero Emission Transportation Association, an electric vehicle industry coalition. Unlike the preceding cases, the target organisation here is not associated with climate denial; ZETA actively advocates for EV adoption. This case illustrates that LACK captures the full complexity of influence networks — including actors whose affiliations span both fossil fuel infrastructure and clean energy advocacy — rather than mapping only clear-cut disinformation actors.
% \paragraph{LACK as a selective filter on COP30 attendance.}
% Crossing LACK with the COP30 participant list shows that the resource functions as a curated filter on the climate-policy field rather than as a comprehensive registry of attendees. Of 47{,}308 registered delegates and 19{,}298 distinct organisations, 0.11\% and 2.87\% respectively are present in LACK, concentrated in the corporate, financial, and lobby-aligned segments LACK is designed to track. The match rate varies fivefold across UNFCCC participation categories---from 2.2\% for government bodies to 10.2\% for the Global Climate Action cluster---and across national delegations, from 8.1\% for Australia and 6.5\% for France to 0\% for several small states. This unevenness reflects LACK's design (corporate-focused, anglophone, Western-centric) and makes the COP30 cross-reference a useful external probe of the resource's effective scope. 
% %Presence in LACK identifies an attendee as a member of the broader climate-policy field; absence carries no such inference.
\begin{sloppypar}
\paragraph{Case study: LACK's prominent collectives at COP30.}
We cross-referenced LACK's 1{,}000 most heavily-described \texttt{lack:Collective} entities (ranked by outgoing IRI links) with the COP30 participant list via a simple heuristic label match.
Of 944 distinct delegate matches, roughly 70\% flow through observer accreditation channels, 18\% appear in high-income national delegations (EU institutions, BNP Paribas, METI, the World Economic Forum), and only 43 in lower- or low-income country delegations.
A qualitative look at the top-20 scope shows the LACK's most densely connected entities: the US climate countermovement (Heartland, Cato, ALEC, Heritage, Manhattan Institute) --- of which only three (US Chamber of Commerce, Shell, ExxonMobil) appear at COP30 at all.
The cross-reference surfaces two findings worth investigation, which were invisible to either resource alone: the US climate countermovement does not engage with the UNFCCC track, and observer accreditation rather than national delegation is the dominant channel through which LACK-tracked organisations participate in COP.
This is one example of the broader family of analyses LACK enables when paired with any external source of named climate-policy agents. 
\end{sloppypar}
